\chapter{民谣：一把吉他能装下多少人间}
\label{ch:genre-folk-deep}

\classquote{心之忧矣，我歌且谣。}{《诗经·魏风·园有桃》}{心里有忧愁，我就唱歌、哼谣；不是为表演，是为把话放出去。}

\section{什么是民谣：把一个人放回歌里}
\label{sec:folk-what}

民谣最容易被误会成“穷版流行歌”：一把木吉他、一个主唱、几段和弦，好像谁都能做。真正好听的民谣不是“少放几轨”，而是\textbf{把一个人放回歌里}——一个具体的人，在一个具体的地方，用具体的声音讲一件具体的事。

它的核心听感通常有三层：

\begin{longtable}{p{0.16\textwidth}p{0.36\textwidth}p{0.36\textwidth}}
\toprule
听感 & 听众先听到什么 & 编排上意味着什么 \\
\midrule
\endfirsthead
\toprule
听感 & 听众先听到什么 & 编排上意味着什么 \\
\midrule
\endhead
人声在前 & 像有人在房间里说话，呼吸、停顿、口音都在 & 人声是主奏，吉他是第二人声；不要用伴奏盖住字 \\
纹理清楚 & 木吉他的颗粒、手指擦弦、房间的底噪 & 少轨、近距离、别把中频糊成一片 \\
情绪有起落 & 主歌收着，副歌才敢放；不是整首一个音量 & 副歌靠加人声、加一把琴、换织体，不靠推子 \\
故事可复述 & 听完记得一个地点、一个动作、一个人 & 编排服务叙事，不服务“氛围滤镜” \\
\bottomrule
\end{longtable}

民谣可以很“土”，也可以很“精致”；可以只有一把琴，也可以有弦乐、口琴、手风琴、轻鼓。判断标准不是轨数，而是\textbf{每一层有没有在帮那个人说话}。如果一层声音去掉后，歌只是变轻了，故事没变，它大概可以走；如果去掉后，某句话失去了落点，它就该留下。

\section{来路：从民歌采集到城市里的独唱}
\label{sec:folk-history}

\subsection{1940—1960 年代：民歌复兴与“一个人的歌”}
\label{sec:folk-history-1940s}

现代民谣的来路不止一条。一条是古老的民歌传统：劳动号子、田歌、叙事谣、酒歌，原本就不是在音乐厅里唱的。另一条是 20 世纪中叶的英文民谣复兴：Woody Guthrie 在 1940 年代把流浪、失业、土地写成歌；Pete Seeger 采集、改编、传唱民歌，把“民歌”从博物馆里拉回现场。{\textsuperscript{[1]}}

1950 年代，The Weavers 等组合让民谣进入商业电台；1960 年代，Bob Dylan 把木吉他、口琴和一种“不修饰”的嗓子推到台前。1963 年的《Blowin' in the Wind》不是靠编曲复杂取胜，而是靠一个人、一把琴、一个反复的问题，让听众自己把答案接下去。{\textsuperscript{[2]}}

\subsection{1970 年代：唱作人把私生活写进歌}
\label{sec:folk-history-1970s}

1970 年代，Joni Mitchell、James Taylor、Carole King 等唱作人把“我”放进歌里：爱情、城市、旅行、成长，不再需要宏大叙事。Joni Mitchell 1971 年的专辑《Blue》几乎不用繁复编配，声音、木吉他、Appalachian dulcimer 和少量鼓，把脆弱感留在最前面。{\textsuperscript{[3]}}

这一时期的遗产是“编曲可以做减法”：不是每首歌都要有完整乐队；一件独特的音色（比如 dulcimer 的金属感）可以比十轨伴奏更有辨识度。

\subsection{1990 年代：中国校园民谣}
\label{sec:folk-history-china-1990s}

中国当代民谣的一个关键节点是 1994 年。大地唱片发行《校园民谣 I》，老狼唱《同桌的你》《睡在我上铺的兄弟》，高晓松、沈庆等人把校园里的青春、离别、宿舍和操场写成歌。{\textsuperscript{[4]}} 它让“民谣”从地下磁带、草坪弹唱，变成可以在电台和电视里被听见的主流声音。

校园民谣的成功也留下一个副作用：很多人以为民谣就是“木吉他＋青春怀旧”。其实它只是把个人经验放到了前台；题材可以是校园，也可以是城中村、公路、工厂、海边。

\subsection{2000—2020 年代：独立民谣与全球本地化}
\label{sec:folk-history-2000s}

2000 年代，英美独立民谣把 banjo、手风琴、复古人声和摇滚动态重新组合；Mumford \& Sons 2009 年的《Little Lion Man》用 banjo、立式贝斯、四拍底鼓和多层和声，把民谣做成能站上大舞台的声音。{\textsuperscript{[5]}} 它不是“回到过去”，而是用旧乐器做新的动态。

中国这边，2010 年代独立民谣通过音乐节、Livehouse、选秀和短视频扩散；赵雷 2016 年发行《成都》，2017 年在《歌手》唱到全场第二，把成都的玉林路、小酒馆和一段私人记忆变成大众共情的城市符号。{\textsuperscript{[6]}} 民谣的“本地化”不是加几句方言，而是让一座城市的生活细节成为歌的主角。

\subsection{史料与出处}
\label{sec:folk-sources}

\begin{longtable}{p{0.06\textwidth}p{0.38\textwidth}p{0.46\textwidth}}
\toprule
编号 & 事实 & 出处 \\
\midrule
\endfirsthead
\toprule
编号 & 事实 & 出处 \\
\midrule
\endhead
1 & 1940 年代 Woody Guthrie、Pete Seeger 与美国民歌复兴 & 财新网《永远的抗议歌手》，\url{https://culture.caixin.com/2014-02-26/100643389.html} \\
2 & 1963 年《Blowin' in the Wind》；Dylan 独自弹吉他、吹口琴 & Bob Dylan 官网专辑页，\url{https://www.bobdylan.com/albums/freewheelin-bob-dylan/} \\
3 & 1971 年《Blue》；Joni Mitchell 弹 dulcimer，James Taylor 弹木吉他，Russ Kunkel 打鼓 & Joni Mitchell 官网资料，\url{https://jonimitchell.com/library/view.cfm?id=2545} \\
4 & 1994 年《校园民谣 I》、老狼《同桌的你》 & 新京报《中国民谣音乐史：民谣从何时起变成了流行歌曲？》，\url{https://m.bjnews.com.cn/detail/157267174414834.html} \\
5 & 2009 年《Little Lion Man》；banjo、立式贝斯、四拍底鼓与和声动态 & Sound On Sound, ``Mumford \& Sons: Little Lion Man''，\url{https://www.soundonsound.com/techniques/mumford-sons-little-lion-man} \\
6 & 赵雷《成都》2016 年发行；编曲/制作与乐手署名 & 维基百科《成都 (歌曲)》，\url{https://zh.wikipedia.org/wiki/%E6%88%90%E9%83%BD_(%E6%AD%8C%E6%9B%B2)} \\
7 & 2017 年《歌手》演唱《成都》，全场第二 & 央视网《赵雷和〈成都〉刷爆朋友圈 民谣的春天来了？》，\url{https://news.cctv.com/2017/02/08/ARTIsg0C6fU6tQYOM20KTuGo170208.shtml} \\
\bottomrule
\end{longtable}

\section{配器与编排：吉他不是伴奏墙，是第二个人声}
\label{sec:folk-craft}

\subsection{先把功能写清楚}
\label{sec:folk-instrumentation}

民谣的编制可以只有两轨：人声、木吉他。也可以有六到八轨：人声、双唱、木吉他、第二吉他/banjo、贝斯、轻鼓、口琴/笛子、弦乐。轨道多不等于民谣，关键是每层都在“说话”。

\begin{longtable}{p{0.15\textwidth}p{0.30\textwidth}p{0.27\textwidth}p{0.18\textwidth}}
\toprule
功能层 & 常见轨道/乐器 & 主要任务 & 常见错误 \\
\midrule
\endfirsthead
\toprule
功能层 & 常见轨道/乐器 & 主要任务 & 常见错误 \\
\midrule
\endhead
人声 & 主唱、双唱、低三度和声 & 叙事、呼吸、口音、情绪落点 & 每个字都叠双唱，像复制人 \\
吉他 & 木吉他指弹/扫弦、第二吉他 & 定调、定速度、给句子留缝 & 音型塞满每拍，主唱没地方换气 \\
低频 & 贝斯、立式贝斯、低音鼓 & 守住根音，让副歌有重量 & 主歌就上完整鼓组，副歌反而没对比 \\
色彩 & 口琴、banjo、手风琴、笛子、弦乐 & 给段落边界、给地方气味 & 色彩从头响到尾，变成背景噪音 \\
空间 & 房间声、近距离底噪、混响 & 让人“在一个地方” & 用低保真滤镜冒充真诚 \\
\bottomrule
\end{longtable}

一个可复现的起步是六轨：Vocal（音频）、Guitar（音频）、Guitar Double（音频，只在副歌）、Bass（音频或 MIDI）、Percussion（音频/MIDI，主歌可静音）、Pad/Harmony（MIDI，最后加）。先把人声和一把吉他立住；如果这两轨不能成立，加多少轨都只是遮住问题。

\subsection{织体：指弹、扫弦、Travis picking}
\label{sec:folk-texture}

木吉他有三种最常用的织体：

\begin{itemize}
\item \textbf{扫弦}：适合副歌、适合“把话推出去”；缺点是主歌也扫会太吵，句尾没有落点。
\item \textbf{分解和弦}：适合主歌叙事；缺点是每小节都一样会变机械，需要在句尾换音型或留半拍。
\item \textbf{Travis picking / 交替低音}：拇指守低音，食指/中指拨旋律；适合“一个人在房间里弹”的亲密感，也最容易弹得太满。
\end{itemize}

写的时候先问：这一句要不要让听众听见呼吸？如果要，吉他在句尾停半拍；如果不要，就让吉他继续走，把人声推过去。民谣的“呼吸”不是音量小，而是\textbf{句子之间有地方换气}。

\subsection{和声：简单的循环，清楚的落点}
\label{sec:folk-harmony}

I–IV–V–vi 是民谣最常用的语法，但简单不等于随便。三个可操作的原则：

\begin{enumerate}
\item \textbf{低音要有方向}：不要每小节都回到主和弦；用转位、挂留、开放和弦让低音线走两步。
\item \textbf{副歌只改一件事}：可以把 vi 换成 I，把大三度换成小三度，或把最后一拍拉长；不要同时换和弦、换速度、换织体。
\item \textbf{开放调弦是为音色，不是为炫技}：DADGAD、Open D 这类调弦会让琴有更长的共鸣和更低的根音；使用时先确认主唱的音域，不要为了调弦把歌唱到不舒服的调。
\end{enumerate}

\subsection{动态：副歌不是推子，是“多一个人”}
\label{sec:folk-dynamic}

民谣最有效的副歌抬升，通常不是把音量推大，而是“多一个人/多一件乐器”：

\begin{longtable}{p{0.18\textwidth}p{0.36\textwidth}p{0.38\textwidth}}
\toprule
段落 & 可以加什么 & 不要做什么 \\
\midrule
\endfirsthead
\toprule
段落 & 可以加什么 & 不要做什么 \\
\midrule
\endhead
主歌 & 人声＋木吉他；鼓可只用刷子/铃鼓 & 不要上完整鼓组；不要把吉他扫满 \\
预副歌 & 贝斯进；第二吉他进；和声只叠句尾 & 不要把和声从头铺到尾 \\
副歌 & 双唱/低三度；轻底鼓；口琴或笛子一句 & 不要把所有乐器同时推到最大 \\
桥段 & 撤掉鼓/贝斯，只留人声＋吉他 & 不要为了“变化”加电子 Drop \\
尾奏 & 回到主歌动机；留白 & 不要用长混响掩盖没有结尾 \\
\bottomrule
\end{longtable}

\subsection{夜半骨架：六轨起步}
\label{sec:folk-yeban}

\begin{longtable}{p{0.13\textwidth}p{0.20\textwidth}p{0.55\textwidth}}
\toprule
轨道 & 类型 & 先写什么 \\
\midrule
\endfirsthead
\toprule
轨道 & 类型 & 先写什么 \\
\midrule
\endhead
Vocal & 音频 & 八小节主歌，句尾留半拍到一拍；先不修音，先确认“像说话” \\
Guitar & 音频 & 分解和弦或 Travis picking；每句结尾留缝，不要每拍都塞 \\
Guitar Double & 音频 & 只在副歌进；音量低于主吉他 3--6 dB，别做成两把主奏 \\
Bass & 音频/MIDI & 主歌只守根音；副歌加入经过音；不要和低音鼓抢同一拍 \\
Percussion & 音频/MIDI & 主歌可静音；副歌用刷子/轻底鼓；不要完整鼓组 \\
Pad/Harmony & MIDI & 最后加；只做长音或句尾和声；不要盖住木吉他中频 \\
\bottomrule
\end{longtable}

顺序：人声 → 吉他 → 贝斯 → 副歌双唱 → 轻鼓 → 色彩。每加一轨，Solo 一次，问自己：“它在帮哪句话？”如果答不上来，就撤。

\section{四首歌，四种“少”的办法}
\label{sec:folk-cases}

\subsection{《Blowin' in the Wind》：一个人就是一支乐队}
\label{sec:folk-case-dylan}

1963 年，Bob Dylan 在《The Freewheelin' Bob Dylan》里唱《Blowin' in the Wind》。官方资料写得很清楚：这首歌里 Dylan 独自用木吉他和口琴伴奏，没有鼓，没有乐队。{\textsuperscript{[2]}} 歌用反复的问题推进，每一段都把“答案”往后推，最后把判断交给听众。

它的编曲智慧是“不解决”：吉他只做稳定的循环，口琴在句尾像第二个说话的人；人声不追求漂亮，追求“这句话非说不可”。整首歌没有一处靠编曲炫技，却能让一整个时代跟着唱。

\textbf{可以偷的一招}：如果你只有一把琴，不要急着加伴奏；先写一个能反复的循环，再决定口琴/第二人声在哪个句尾出现。一个人的歌，留白就是乐队。

\subsection{《A Case of You》：音色比轨数更会说话}
\label{sec:folk-case-joni}

Joni Mitchell 1971 年的《A Case of You》收录在《Blue》。官网资料显示，这首歌的伴奏很克制：Joni 弹 Appalachian dulcimer，James Taylor 弹木吉他，Russ Kunkel 打鼓；声音被放在最前面。{\textsuperscript{[3]}} dulcimer 的金属弦音色和木吉他的暖颗粒并置，像两种不同的“我”在对话。

它没有把爱情写成大合唱，而是让音色承担情绪：dulcimer 的亮、吉他的软、鼓的轻，三件乐器各说一句。副歌不靠加轨，靠人声从低语到高音的滑动。

\textbf{可以偷的一招}：给一首歌找一件“不属于标准乐队”的音色——dulcimer、口琴、手风琴、甚至一把调得很低的木吉他。让它只出现两三次，比从头铺到尾更有记忆点。

\subsection{《Little Lion Man》：用“进人”做副歌}
\label{sec:folk-case-mumford}

Mumford \& Sons 2009 年的《Little Lion Man》是他们的首支单曲，收录在《Sigh No More》。Sound On Sound 的制作访谈提到：专辑大部分是乐队在录音室现场演奏，之后再加叠鼓；鼓主要是四拍底鼓和地板鼓，偶尔加军鼓；编制里有立式贝斯、banjo、钢琴、Marcus Mumford 的人声与木吉他，以及大量和声。{\textsuperscript{[5]}}

它的副歌不是“音量变大”，而是“人越来越多”：先是木吉他，接着 banjo、钢琴、贝斯、底鼓一层层进来，最后和声把整个房间抬起来。banjo 的颗粒很抢，但制作上刻意让它“切得出来又不变成 banjo 专场”，这就是动态分配。

\textbf{可以偷的一招}：副歌前不要把所有乐器都放进去。让每件乐器在副歌的不同拍点进入，听众会感觉“歌在长”，而不是“推子被推大”。

\subsection{赵雷《成都》：让一座城市成为第二主唱}
\label{sec:folk-case-chengdu}

赵雷 2016 年发行《成都》，2016 年 12 月收录进专辑《无法长大》。制作署名里，赵雷与喜子负责编曲/制作，张岭弹贝斯，贝贝打鼓，柳森弹钢琴，祝子吹笛子，赵雷与喜子弹箱琴，另有弦乐编写。{\textsuperscript{[6]}} 2017 年他在《歌手》唱这首歌，拿到全场第二，把独立民谣推到大众面前。{\textsuperscript{[7]}}

这首歌的编配没有把成都做成旅游宣传片：主歌是木吉他与人声的近景，贝斯和鼓很轻，笛子和弦乐只在段落边界出现，像街角突然传来的一声。它让“玉林路”“小酒馆”这些具体地名承担叙事，而不是靠形容词堆情绪。

\textbf{可以偷的一招}：写城市不要写“这座城市好美”，写三个只有本地人知道的物件/地点/时间；编配上只给其中一个物件一件音色（比如笛子、手风琴、一段环境声），不要全上。

\section{三个练习：在夜半把民谣拆开}
\label{sec:folk-practice}

\begin{exercise}
\textbf{同一段词，三种吉他织体。} 用八小节主歌，做 A 分解和弦、B 扫弦、C Travis picking 三个版本。每个版本都录人声，盲听比较：哪个最像“说话”，哪个最像“表演”，哪个副歌最好接。A 亲密但容易平；B 有推力但容易吵；C 有纹理但容易弹满。没有正确答案，只有取舍。
\end{exercise}

\begin{exercise}
\textbf{撤轨测试。} 先做一版“完整乐队”Demo（人声、吉他、贝斯、鼓、Pad、和声），再撤到只剩人声＋吉他，然后一次只加回一件乐器。每次加回都问：它帮了哪句话？如果只是让歌变“满”，就不要加。最后留下三到五轨，通常比八轨更像民谣。
\end{exercise}

\begin{exercise}
\textbf{一座城市，三个具体物件。} 写八句词，不许用“美丽”“孤独”“远方”这类大词；只能写三个具体地点/物件/时间（例如一家店、一班车、一场雨）。配一条 I–IV–V–vi 循环，主歌只用吉他，副歌加一件色彩乐器（笛子/口琴/手风琴）。先保证清唱能站住，再谈编配。
\end{exercise}

\section{章末扩展阅读}
\label{sec:folk-reading}

\begin{sloppypar}
\begin{itemize}
\item \textbf{Bob Dylan 官网《The Freewheelin' Bob Dylan》专辑页}：\url{https://www.bobdylan.com/albums/freewheelin-bob-dylan/}。核对《Blowin' in the Wind》的录制、吉他/口琴编制与专辑语境。
\item \textbf{Joni Mitchell 官网《Blue》资料}：\url{https://jonimitchell.com/music/album.cfm?id=5}。看《A Case of You》在专辑里的位置，以及 dulcimer、木吉他、鼓的克制编配。
\item \textbf{Sound On Sound《Mumford \& Sons: Little Lion Man》}：\url{https://www.soundonsound.com/techniques/mumford-sons-little-lion-man}。从制作访谈理解 banjo、立式贝斯、四拍底鼓与和声动态。
\item \textbf{央视网《赵雷和〈成都〉刷爆朋友圈》}：\url{https://news.cctv.com/2017/02/08/ARTIsg0C6fU6tQYOM20KTuGo170208.shtml}。核对《成都》在《歌手》后的传播与独立民谣主流化。
\end{itemize}
\end{sloppypar}

本章扩展阅读链接已于 2026 年 11 月 27 日逐条核对，均可访问。若遇地区限制，可在对应平台的官方账号内搜索标题。

\section{本章要点}
\label{sec:folk-takeaways}

\begin{itemize}
\item 民谣不是“穷版流行歌”，是把一个具体的人、地点、时刻放回歌里。
\item 人声是主奏，吉他是第二个人声；吉他的每一拍都要给人声留换气的地方。
\item 副歌的抬升优先靠“多一个人/多一件乐器”，不要先推音量。
\item 音色比轨数更会说话：一件不属于标准乐队的音色，只出现两三次，就够记住。
\item 城市叙事靠具体物件，不靠大词；编配只给其中一个物件一件音色。
\item 先让清唱和一把吉他站住；站不住，加轨只是遮住问题。
\end{itemize}

\begin{helpmapping}
\begin{longtable}{p{0.45\textwidth}p{0.45\textwidth}}
\toprule
小节 & 帮助锚点 \\
\midrule
什么是民谣 & \texttt{arrangement-genre-folk-what} \\
来路 & \texttt{arrangement-genre-folk-history} \\
配器与编排 & \texttt{arrangement-genre-folk-craft} \\
四首歌 & \texttt{arrangement-genre-folk-cases} \\
夜半骨架 & \texttt{arrangement-genre-folk-yeban} \\
练习 & \texttt{arrangement-genre-folk-practice} \\
扩展阅读 & \texttt{arrangement-genre-folk-reading} \\
\bottomrule
\end{longtable}
\end{helpmapping}
